\section{Three bands from an actual ideal selection}
\label{sec:pencil-ideal-kernels}
\index{Bergman kernel!pencil selection}

For an isolated logarithmic pole of coefficient one on a complex surface, the degree-$m$ multiplier ideal permits vanishing of order $m-1$, while ordinary logarithmic domination requires order $m$. At the two base points of the pencil below, these conditions produce three neighboring monomial bands. All three contribute to the kernel measure, although the full ideal and ordinary competitor classes agree.

Let $X=\mathbb P^1\times\mathbb P^1$, $L=\mathcal O(1,1)$ with its product Fubini--Study metric $h$, and $\theta=c_1(L,h)$. Write $z_j=Z_j/T_j$ and set
\[
 \psi=\max\{\log|T_1T_2|_h^2,\log|Z_1Z_2|_h^2\},
 \qquad p_+=(0,\infty),\quad p_-=(\infty,0).
\]
The only poles are $p_\pm$. The local frame potentials are maxima of logarithms of squared holomorphic functions, so $\psi$ is genuinely $\theta$-psh. On the torus put
\[
 x_j=\log|z_j|^2,\qquad s=x_1+x_2,\qquad
 \varphi_X=\log(1+e^{x_1})+\log(1+e^{x_2}).
\]
Then $\psi=\max(0,s)-\varphi_X$.

For this model the norm is taken on a specified compact set $K\Subset(\mathbb C^*)^2$, with $x\in[-1,2]^2$ and all angles. It has nonempty interior. Let
\[
 w=\frac{x_1^2+x_2^2}{2}-\varphi_X\quad\hbox{on }K,
 \qquad d\nu=\frac19\,dx_1\,dx_2\,d\tau_1\,d\tau_2,
\]
where each $d\tau_j$ is angular Haar probability. Thus $\nu$ is a probability supported on $K$, rather than a smooth positive volume on all $X$. Set
\[
 \begin{aligned}
 V_m&=H^0(X,L^m\otimes\mathcal I(m\psi)),
 &N_m&=\dim V_m,\\
 \|\sigma\|_m^2&=\int_K|\sigma|_{h^m}^2e^{-mw}\,d\nu,
 &\mathcal B_m&=\sum_j|\sigma_{m,j}|_{h^m}^2,
 \end{aligned}
\]
for an orthonormal basis. The norm is positive definite by the identity theorem on the interior of $K$. Use $U_m=m^{-1}\log\mathcal B_m$ on $X$ and $B_m=e^{-mw}\mathcal B_m$ on $K$, in the convention of \cref{sec:partial-selected-space}. Only the specified values of $w$ enter the norm and envelope.

\Needspace{12\baselineskip}
\begin{lemma}\label{lem:pencil-ideal-ordinary-class}
For every genuine $v\in\PSH(X,\theta)$,
\[
 v\preceq_{\mathcal I}\psi
 \quad\Longleftrightarrow\quad v\leq\psi+C_v
 \quad\hbox{for some finite }C_v.
\]
The ideal and ordinary envelopes therefore agree for every common obstacle. At either pole, for every real $t>0$,
\[
 \mathcal I(t\psi)=\mathfrak m_{p_\pm}^{\max\{\lfloor t\rfloor-1,0\}}.
\]
Away from the poles the ideal is the full structure sheaf.
\end{lemma}
\begin{proof}
Near $p_+$ use $\zeta_1=z_1$, $\zeta_2=1/z_2$ and the frame $T_1Z_2$. The two pencil sections have coefficients $\zeta_2,\zeta_1$, so
\[
 \psi=\log\max(|\zeta_1|^2,|\zeta_2|^2)
       -\log(1+|\zeta_1|^2)-\log(1+|\zeta_2|^2)
      =\log|\zeta|^2+O(1).
\]
The same calculation applies at $p_-$. The ideal formula is \cref{lem:smooth-center-multiplier-ideal} with normal dimension two and coefficient one, including the strict integrability endpoint.

Ideal domination gives $\nu(v,p_\pm)\geq1$ by \cref{thm:local-i-inputs}. For a genuine local psh frame potential $\phi$, consider its radial maximum
\[
 M(a)=\sup_{|\zeta|\leq e^{a/2}}\phi(\zeta)
\]
inside a fixed coordinate ball. It is finite and convex: the one-variable three-circles maximum on each complex line is convex, and their supremum gives $M$. Identically $-\infty$ restrictions do not affect this supremum. Its limiting left slope is the Lelong number in \cref{sec:analytic-types}. For $a<a_0$, convexity and this slope give
\[
 M(a)\leq M(a_0)+\nu(v,p_\pm)(a-a_0)\leq a+C.
\]
Consequently $v\leq\psi+C$ near both poles. The same upper bound, after increasing $C$, holds on the compact complement because $v$ is bounded above and $\psi$ is finite and continuous there. Conversely ordinary domination implies all-weight ideal inclusion by integrability. The competitor classes coincide before imposing any obstacle.
\end{proof}

The pencil also describes competitors which have no symmetry.

\Needspace{13\baselineskip}
\begin{lemma}\label{lem:pencil-full-competitors}
On $X\setminus\{p_+,p_-\}$ set
\[
 f=[T_1T_2:Z_1Z_2],\qquad
 \psi_{\rm FS}=\log(|T_1T_2|_h^2+|Z_1Z_2|_h^2).
\]
Every member of the class in \cref{lem:pencil-ideal-ordinary-class} has the form
\[
 v=\psi_{\rm FS}+u\circ f,
 \qquad u\in\PSH(\mathbb P^1,\omega_{\rm FS}),
\]
with its canonical extension at the two poles. Conversely every such $u$ gives a member of that class.
\end{lemma}
\begin{proof}
For $\xi\in\mathbb C^*$, the pencil fiber
\[
 C_\xi=\{Z_1Z_2=\xi T_1T_2\}
\]
is a smooth rational curve of bidegree $(1,1)$, containing both poles. Its $\theta$-degree is two. Restriction of $v$ to this connected curve is either genuine or identically $-\infty$. In the genuine case its positive curvature measure has total mass two by Stokes. The local upper bounds from \cref{lem:pencil-ideal-ordinary-class} give an atom of mass at least one at each pole, using the planar atom identity in \cref{lem:planar-multiplier-ideal}. These atoms exhaust the measure. The curvature of $\psi|_{C_\xi}$ consists of exactly the same two unit atoms. Thus $v-\psi$ is harmonic in distributions on the compact curve and is constant.

On the torus use holomorphic coordinates $(z_1,\xi=z_1z_2)$. Local integrability and Fubini show that almost every fiber restriction is genuine. Each good fiber is finite off its poles after the preceding constant-difference conclusion. Hence the restriction to the transverse section $S=\{z_1=1\}\simeq\mathbb P^1$ is genuine. Put $u=v|_S+\log2$. Since $\theta|_S=\omega_{\rm FS}$ and $\psi_{\rm FS}|_S=-\log2$, the fiber constants give the displayed identity almost everywhere on the torus. Both sides have genuine local psh frame potentials, so equality of canonical representatives gives equality everywhere off the poles. Their logarithmic upper bounds give value $-\infty$ at the poles.

Conversely $\theta+\ddc(\psi_{\rm FS}+u\circ f)=f^*(\omega_{\rm FS}+\ddc u)\geq0$ off the poles. In a pole chart the frame potential is bounded above by $\log(|\zeta_1|^2+|\zeta_2|^2)+C$. Its canonical removable psh extension is therefore global and has the required ordinary upper type.
\end{proof}

\Needspace{17\baselineskip}
\begin{theorem}\label{thm:pencil-ideal-envelope}
For the specified $K,w$, the full ideal envelope of \eqref{eq:hermitian-partial-i-envelope} is
\[
 \begin{aligned}
 E_{\theta,K}[\psi]_{\mathcal I}(w)&=P=g(s)-\varphi_X,\\
 g(s)&=\sup_{0\leq t\leq1}(ts-t^2)
      =\begin{cases}
         0,&s\leq0,\\
         s^2/4,&0\leq s\leq2,\\
         s-1,&s\geq2.
        \end{cases}
 \end{aligned}
\]
Here $P$ has its canonical global extension, with $P(p_\pm)=-\infty$ and $\psi-1\leq P\leq\psi$.
\end{theorem}
\begin{proof}
For an arbitrary competitor, \cref{lem:pencil-full-competitors} gives on the affine base
\[
 \chi(\xi)=\log(1+|\xi|^2)+u(\xi),
 \qquad v=\chi(z_1z_2)-\varphi_X.
\]
On the interior of $K$, the quasi-everywhere obstacle becomes an everywhere inequality by the submean inequality and continuity of $w+\varphi_X$. Evaluating at $x_1=x_2=s/2$ gives
\[
 \chi(\xi)\leq s^2/4\quad(-2<s<4)
\]
for every argument of $\xi$. At $s=0$ this bounds $\chi$ by zero on $|\xi|=1$. The subharmonic maximum principle bounds it by zero on the entire disk. At infinity use the genuine local potential
\[
 \chi_\infty(\eta)=\chi(1/\eta)+\log|\eta|^2.
\]
The obstacle at $s=2$ bounds it by $-1$ on $|\eta|=e^{-1}$. The same maximum principle gives $\chi(\xi)\leq s-1$ for $s\geq2$. Together with the obstacle for $0\leq s\leq2$, this proves $\chi\leq g$ without invariance of the competitor.

Conversely $g$ is convex with slopes in $[0,1]$ and defines a genuine base Fubini--Study psh potential $g(\log|\xi|^2)-\log(1+|\xi|^2)$. The preceding pencil lemma gives the global candidate $P$. Directly
\[
 \max(0,s)-1\leq g(s)\leq\max(0,s),\qquad
 g(s)\leq s^2/4\leq(x_1^2+x_2^2)/2.
\]
Thus $P$ belongs to the full ideal class and satisfies the obstacle everywhere on $K$. It dominates every competitor, proving the envelope identity including its canonical values.
\end{proof}

\Needspace{13\baselineskip}
\begin{proposition}\label{prop:pencil-three-band-space}
For every $m\geq1$ the entire selected space is
\[
 V_m=\operatorname{span}\{Z_1^aT_1^{m-a}Z_2^bT_2^{m-b}:
              0\leq a,b\leq m,\ |a-b|\leq1\},
 \qquad N_m=3m+1.
\]
Its Hilbert growth exponent is one. These spaces are not multiplicative. Ordinary upper type for a degree-$m$ logarithmic section would select only $a=b$. The measure $\nu$ is Bernstein--Markov for the specified $(K,w)$.
\end{proposition}
\begin{proof}
At $p_+$ the monomial $(a,b)$ has local exponents $(a,m-b)$, and at $p_-$ it has $(m-a,b)$. Both total orders must be at least $m-1$ by \cref{lem:pencil-ideal-ordinary-class}. Distinct exponent pairs cannot cancel, so these conditions select exactly $|a-b|\leq1$. The three bands have $m+1,m,m$ terms. Ordinary upper type would require total order $m$ at both poles and leave only the diagonal. The product of two offset-one sections has offset two and lies outside the selected space in the sum degree; such products fail in every pair of positive degrees. The Hilbert growth exponent is consequently not an assertion about the Iitaka dimension of a graded section ring.

Angular integration makes all global monomials orthogonal, including those outside the selected space. With $t\in[0,1]$ set
\[
 I_m(t)=\int_{-1}^2e^{m(tx-x^2/2)}dx
       =e^{mt^2/2}J_m(t),\qquad
 J_m(t)=\int_{-1}^2e^{-m(x-t)^2/2}dx.
\]
The exact squared norm is
\begin{equation}\label{eq:pencil-gaussian-norm}
 A_m(a,b)=\frac19 I_m(a/m)I_m(b/m).
\end{equation}
For $c=\int_{-1}^1e^{-y^2/2}dy$, uniformly in $t\in[0,1]$,
\[
 \frac c{\sqrt m}\leq J_m(t)\leq\sqrt{\frac{2\pi}m},\qquad
 J_m(t)=\sqrt{\frac{2\pi}m}
                  \bigl(1+O(m^{-1/2}e^{-m/2})\bigr).
\]
The lower bound integrates over a centered interval of radius $m^{-1/2}$; each omitted Gaussian tail for the last assertion is at most $e^{-m/2}/m$. Hence
\begin{equation}\label{eq:pencil-gaussian-norm-bounds}
 \frac{c^2}{9m}e^{(a^2+b^2)/(2m)}
 \leq A_m(a,b)\leq\frac{2\pi}{9m}e^{(a^2+b^2)/(2m)}.
\end{equation}
For $x\in[-1,2]^2$, completing squares bounds every weighted orthonormal contribution by $9m/c^2$. Summing all $(m+1)^2$ global monomials and using the evaluation inequality gives
\[
 \sup_K|\sigma|_{h^m}e^{-mw/2}
       \leq C m^{3/2}\|\sigma\|_m
       \quad(\sigma\in H^0(X,L^m)).
\]
This polynomial bound proves Bernstein--Markov. On the selected space the same argument gives $Cm$.
\end{proof}

\Needspace{11\baselineskip}
\begin{theorem}\label{thm:pencil-ideal-kernel-limits}
The kernels for the specified norm satisfy, with extended-real interpretation at the poles,
\begin{equation}\label{eq:pencil-global-lower}
 U_m\geq P+\frac{\log(9m/(2\pi))}{m}-\frac1{4m^2}.
\end{equation}
They converge locally uniformly on $X\setminus\{p_+,p_-\}$, in every finite $L^p$ against a smooth volume on $X$ at rate $O(\log m/m)$, and extended-pointwise on $X$. Their closed curvature currents converge weakly. Moreover $e^{U_m}\to e^P$ uniformly on $X$. For each fixed $m\geq2$, the additive error $U_m-P$ is unbounded above near either pole.

Define a probability $\mu$ by
\[
 \int_X f\,d\mu=\int_0^1\int f(e^{t/2}\tau_1,e^{t/2}\tau_2)
                          \,d\tau_1\,d\tau_2\,dt.
\]
For every continuous test function on $X$,
\[
 N_m^{-1}B_m\nu\rightharpoonup\mu,\qquad
 m^{-1}B_m\nu\rightharpoonup3\mu.
\]
The ambient normalization $2m^{-2}B_m\nu$ tends to zero in total variation. The envelope has $\langle(\theta+\ddc P)^2\rangle=0$.
\end{theorem}
\begin{proof}
The diagonal sections belong to the actual space. Their upper norm bound in \eqref{eq:pencil-gaussian-norm-bounds} gives
\[
 \mathcal B_m\geq\frac{9m}{2\pi}e^{-m\varphi_X}
                     \max_{0\leq j\leq m}e^{js-j^2/m}.
\]
For $0\leq s\leq2$ the nearest grid point to $s/2$ loses at most $1/(4m^2)$ in $g(s)$. Outside this interval the maximizing endpoint is on the grid. This proves \eqref{eq:pencil-global-lower} on the torus.

For an offset-one term $(a,b)=(j+1,j)$, its exponent after the norm estimate is at most both
\[
 js-j^2/m+x_1,\qquad
 (j+1)s-(j+1)^2/m+1-x_2.
\]
The other band interchanges the indices. Put
\[
 R(x)=1+\min(e^{x_1},e^{1-x_2})
           +\min(e^{x_2},e^{1-x_1}).
\]
The lower norm estimate and the count give
\begin{equation}\label{eq:pencil-global-upper}
 U_m\leq P+\frac{C+2\log(m+1)+\log R(x)}m.
\end{equation}
The function $R$ extends continuously across coordinate divisors away from the poles. In a chart at $p_+$ put $A=|\zeta_1|^2$, $B=|\zeta_2|^2$. The possibly large term satisfies
\[
 \min(1/B,e/A)\leq(e+1)/(A+B);
\]
the other term is bounded. The same holds at $p_-$. Thus $R\leq C(1+r^{-2})$ near either pole, for the local radius $r=|\zeta|$, and $\log R$ belongs to every finite $L^p$ against smooth volume.

For $m\geq2$ all selected sections vanish at the poles. The two diagonal endpoint sections have exactly these common zeros, so there are no others. Away from the poles $P$ and $U_m$ are continuous, so both inequalities extend from the torus. They give the asserted local uniform and $L^p$ convergence. Distributional differentiation of the $L^1$ limit gives weak convergence of closed currents. At the poles both potentials are $-\infty$ for $m\geq2$, which also proves the global extended-pointwise statement.

For a fixed $m\geq2$, the lowest local terms include every binary monomial of degree $m-1$. Their positive squared sum is comparable to $r^{2(m-1)}$ in all directions. Hence
\[
 U_m=(1-1/m)\log r^2+O_m(1),\qquad P=\log r^2+O(1).
\]
The additive error therefore tends to $+\infty$ near each pole. Nevertheless \eqref{eq:pencil-global-upper} gives, uniformly in $m\geq2$ and $r<1$,
\[
 e^{U_m}\leq Cr^{2-2/m}\leq Cr,\qquad e^P\leq Cr^2.
\]
Combined with uniform convergence on compact subsets off the poles, these bounds give uniform exponential convergence.

Each basis contribution to $B_m\nu$ is exactly the product of truncated Gaussian probabilities centered at $(a/m,b/m)$, with independent angular Haar probability. Completing squares in \eqref{eq:pencil-gaussian-norm} and using its uniform bounds prove concentration at these centers, uniformly over all three bands. For example, the probability of $|x-t|\geq\eta$ is at most $C_\eta\sqrt m\,e^{-m\eta^2/2}$. Uniform continuity of a test function then reduces the sum to the three lattice Riemann sums. Each band has asymptotic proportion $1/3$ in $N_m$, and each center sequence tends to the same diagonal parameter $t\in[0,1]$. This proves both positive limits. The ambient-normalized mass is $2(3m+1)/m^2\to0$, proving total-variation convergence to zero.

Off the poles the curvature is the pullback of the base current defined by $g$. The function $g$ is $C^{1,1}$; its local complex Hessian has rank at most one. Smooth convex approximations and the locally bounded product identify its square with zero, as in the Hessian argument of \cref{sec:newton-selected-kernels}. The nonpluripolar product assigns no mass to the two poles, so its global square is zero. The positive-top-mass hypothesis of \cref{thm:fujiki-partial-kernel-limit} is therefore absent in this model.
\end{proof}

Parametrize the real three-dimensional shell $S$ by $\Psi(t,\tau_1,\tau_2)=(e^{t/2}\tau_1,e^{t/2}\tau_2)$, for $0<t<1$. For a smooth positive Hermitian form $\omega$, put $\Psi^*d\operatorname{vol}_{\omega,S}=J_\omega\,dt\,d\tau_1\,d\tau_2$, as in \cref{sec:newton-selected-kernels}. The densities of $\mu$ and $3\mu$ are $1/J_\omega$ and $3/J_\omega$; endpoints have zero mass. Closedness and invariance of $\omega$ are unnecessary. The mixed measure $\langle\theta+\ddc P\rangle\wedge\omega$ has a locally bounded density off the poles, gives $S$ zero mass and cannot reproduce the positive kernel limit.
